\section{Results}

The results below are populated after running the experiment pipeline (Phase 3). Refer to Table~\ref{tab:correlation} for ranking correlations, Table~\ref{tab:precision} for precision@N, and Table~\ref{tab:ablation} for the threshold ablation.

\subsection{Correlation analysis}

TBD. We expect the weighted h-index to correlate strongly with the conventional h-index for top researchers (both agree on stars) but to diverge in the middle of the ranking distribution, where citation-quality differences matter most.

\subsection{Precision@N}

TBD. We expect the weighted h-index to achieve higher precision@N than the conventional h-index at identifying influential researchers, mirroring Zhu et al.'s finding that the hip-index (AveP 14\%) outperformed the conventional h-index (AveP 10\%) at identifying ACL Fellows.

\subsection{Threshold ablation}

TBD. Higher relevance thresholds ($\tau \in \{0.5, 0.7\}$) should retain fewer but more meaningful citations, potentially improving discrimination among mid-tier researchers while preserving the top ranking.

\subsection{Discussion}

TBD. Interpret results in light of Zhu et al. (2015) and discuss limitations: abstract-only relevance scoring, LLM variability, OpenAlex coverage, and cost of LLM calls at scale.
